% from ega0-1-fr.tex
\oldpage[0\textsubscript{I}]{18}
qui, à l'élément $(n'/s')\otimes(a/s)$, fait correspondre l'élément
$(n'\otimes a)/(\varphi(s')s)$; on vérifie en effet séparément que lorsqu'on
remplace $n'/s'$ (resp. $a/s$) par une autre expression du même élément,
$(n'\otimes a)/(\varphi(s')s)$ ne change pas; d'autre part, on peut définir un
homomorphisme réciproque de $\tau$ en faisant correspondre à $(n'\otimes a)/s$
l'élément $(n'/1)\otimes(a/s)$: on utilise le fait que
$S^{-1}(N'\otimes_{A'}A_{[\varphi]})$ est canoniquement isomorphe à
$(N'\otimes_{A'}A_{[\varphi]})\otimes_A S^{-1}A$ (1.2.5), donc aussi à
$N'\otimes_{A'}(S^{-1}A)_{[\psi]}$, en désignant par $\psi$ l'homomorphisme
composé $a'\mapsto\varphi(a')/1$ de $A'$ dans $S^{-1}A$.
\end{env}

\begin{env}[1.5.5]
\label{0.1.5.5-fr}
Si $M'$ et $N'$ sont deux $A'$-modules, en composant les isomorphismes (1.3.4)
et (1.5.4), on obtient un isomorphisme
\[
  S'^{-1}M'\otimes_{S'^{-1}A'}S'^{-1}N'
  \otimes_{S'^{-1}A'}S^{-1}A
  \xrightarrow{\sim}
  S^{-1}(M'\otimes_{A'}N'\otimes_{A'}A).
\]
De même, si $M'$ admet une présentation finie, on a en vertu de (1.3.5) et
(1.5.4), un isomorphisme
\[
  \operatorname{Hom}_{S'^{-1}A'}(S'^{-1}M',S'^{-1}N')
  \otimes_{S'^{-1}A'}S^{-1}A
  \xrightarrow{\sim}
  S^{-1}(\operatorname{Hom}_{A'}(M',N')\otimes_{A'}A).
\]
\end{env}

\begin{env}[1.5.6]
\label{0.1.5.6-fr}
Sous les hypothèses de (1.5.1), soit $T$ (resp. $T'$) une seconde partie
multiplicative de $A$ (resp. $A'$) telle que $S\subset T$ (resp. $S'\subset T'$)
et $\varphi(T')\subset T$. Alors le diagramme
\[
\begin{tikzcd}[column sep=huge,row sep=large]
S'^{-1}A' \arrow[r,"\varphi^{S'}"] \arrow[d,"\rho^{T',S'}"']
  & S^{-1}A \arrow[d,"\rho^{T,S}"] \\
T'^{-1}A' \arrow[r,"\varphi^{T'}"']
  & T^{-1}A
\end{tikzcd}
\]
est commutatif. Si $M$ est un $A$-module, le diagramme
\[
\begin{tikzcd}[column sep=huge,row sep=large]
S'^{-1}(M_{[\varphi]}) \arrow[r,"\sigma"] \arrow[d,"\rho^{T',S'}"']
  & (S^{-1}M)_{[\varphi^{S'}]} \arrow[d,"\rho^{T,S}"] \\
T'^{-1}(M_{[\varphi]}) \arrow[r,"\sigma"']
  & (T^{-1}M)_{[\varphi^{T'}]}
\end{tikzcd}
\]
est commutatif. Enfin, si $N'$ est un $A'$-module, le diagramme
\[
\begin{tikzcd}[column sep=huge,row sep=large]
(S'^{-1}N')\otimes_{S'^{-1}A'}(S^{-1}A)_{[\varphi^{S'}]}
  \arrow[r,"\tau","{\sim}"'] \arrow[d]
  & S^{-1}(N'\otimes_{A'}A_{[\varphi]}) \arrow[d,"\rho^{T,S}"] \\
(T'^{-1}N')\otimes_{T'^{-1}A'}(T^{-1}A)_{[\varphi^{T'}]}
  \arrow[r,"\tau"',"{\sim}"]
  & T^{-1}(N'\otimes_{A'}A_{[\varphi]})
\end{tikzcd}
\]
est commutatif, la flèche verticale de gauche étant obtenue en appliquant
$\rho_{N'}^{T',S'}$ à $S'^{-1}N'$ et $\rho_A^{T,S}$ à $S^{-1}A$.
\end{env}

